\section{Complemento de model-based RL: generación de datos sintéticos}

La primera mitad de esta clase cierra primero lo que quedó pendiente de model-based RL en la clase anterior. El punto de partida de Chelsea Finn es muy claro: si ya aprendimos un modelo capaz de predecir estados futuros, además de usarlo para planning, ¿podemos usarlo también para \textbf{producir más datos de entrenamiento} y, a su vez, entrenar una política más fuerte y más barata?

\begin{figure}[H]
\centering
\includegraphics[width=0.9\textwidth]{slides-images/slide-04.jpg}
\caption{Planning with learned models: planificar, recolectar datos, actualizar el modelo y replanificar periódicamente}
\end{figure}

\subsection{De planning a policy optimization}

La clase anterior trató de «dado el estado actual, usar un learned model para buscar en tiempo de prueba secuencias de acciones futuras». Esta clase va un paso más allá: si buscar repetidamente en tiempo de prueba resulta demasiado costoso, ¿es posible convertir la capacidad predictiva que aporta el modelo en una herramienta de aumento de datos para la etapa de aprendizaje de la política?

\begin{importantbox}{El problema que esta sección busca resolver en realidad}
Planning y policy optimization cumplen funciones distintas:
\begin{itemize}
    \item La ventaja de planning es su flexibilidad: si cambia el objetivo, se puede volver a calcular.
    \item La desventaja de planning es que su test-time compute es alto y, con un horizon largo, la acumulación del error del modelo es notable.
    \item La ventaja de aprender directamente la política es que su despliegue es barato.
    \item La desventaja de aprender directamente la política es que, si los datos de entrenamiento no bastan, es difícil aprenderla de forma estable.
\end{itemize}
Por lo tanto, la pregunta más natural es: \textbf{cómo usar un learned model para producir más datos útiles para la política y, al mismo tiempo, evitar que el error del modelo se amplifique hasta volverse incontrolable.}
\end{importantbox}

\subsection{Aumento de datos al estilo Dyna}

\begin{figure}[H]
\centering
\includegraphics[width=0.9\textwidth]{slides-images/slide-05.jpg}
\caption{La decisión central de model-based policy optimization: desde dónde iniciar el rollout y qué tan largo hacerlo}
\end{figure}

El curso plantea este problema de manera muy concreta. Supongamos que ya recolectamos varias trayectorias en el entorno real; entonces, al generar datos con el modelo, hay al menos tres estrategias:
\begin{enumerate}
    \item Hacer rollouts de trayectorias muy largas a partir del estado inicial.
    \item Hacer solo rollouts cortos a partir del estado inicial.
    \item Hacer rollouts cortos a partir de \textbf{todos los estados que ya aparecen} en el conjunto de datos.
\end{enumerate}

La tercera estrategia es la práctica central que recomienda esta clase, porque equilibra la cobertura con el control del error.

\begin{figure}[H]
\centering
\includegraphics[width=0.9\textwidth]{slides-images/slide-06.jpg}
\caption{Un término medio más seguro: partir de todos los estados de los datos y generar synthetic rollouts locales}
\end{figure}

\begin{knowledgebox}{Por qué «rollout corto + todos los estados» es la opción más segura}
\begin{enumerate}
    \item Si solo se generan trayectorias largas desde el estado inicial, el error del modelo se amplifica continuamente a lo largo de la cadena.
    \item Si solo se generan trayectorias cortas desde unos pocos puntos de partida, la cobertura de los estados posteriores resulta muy insuficiente.
    \item Hacer rollouts cortos a partir de todos los estados de los datos equivale, en esencia, a extrapolar localmente cerca de la distribución real, lo que mejor equilibra estabilidad y cobertura.
\end{enumerate}
\end{knowledgebox}

\begin{figure}[H]
\centering
\includegraphics[width=0.9\textwidth]{slides-images/slide-07.jpg}
\caption{Algoritmo completo: los datos reales actualizan el modelo y los rollouts cortos del modelo generan muestras de entrenamiento adicionales}
\end{figure}

\subsection{Cuándo vale la pena aprender un modelo del mundo}

Model-based RL resulta muy atractivo porque parece «imaginar el futuro». Sin embargo, el criterio de Chelsea Finn es mesurado: \textbf{lo decisivo no es si el modelo suena sofisticado, sino si en verdad es más fácil de aprender que la política.}

\begin{figure}[H]
\centering
\includegraphics[width=0.9\textwidth]{slides-images/slide-08.jpg}
\caption{Ventajas y costos de model-based RL: equilibrio entre eficiencia de datos, independencia de la tarea y complejidad del entrenamiento}
\end{figure}

Los principios de decisión que da la ponente pueden resumirse en una tabla breve:

\begin{table}[H]
\centering
\begin{tabular}{p{0.2\textwidth} p{0.28\textwidth} p{0.28\textwidth}}
\toprule
Característica del escenario & Más probablemente favorable a model-based & Riesgos \\
\midrule
La interacción real es costosa & Se puede mejorar la eficiencia de datos con synthetic rollouts & La política absorbe el sesgo del modelo \\
La dinámica es relativamente regular & Es más fácil aprender una aproximación útil del modelo del mundo & Es fácil subestimar el error con horizon largo \\
La recompensa cambia con frecuencia & Un mismo modelo puede servir a varios objetivos posteriores & No implica que el modelo mismo optimice el desempeño en la tarea \\
Entorno muy estocástico o parcialmente observable & No necesariamente adecuado & Puede ser más difícil que aprender directamente la política \\
\bottomrule
\end{tabular}
\caption{Cuándo considerar primero model-based RL}
\end{table}

\begin{warningbox}{No hay que tomar model-based como la vía más fuerte por defecto}
El modelo del mundo requiere entrenamiento y ajuste de hiperparámetros propios, y además introduce nuevas fuentes de error. Si la dinámica del entorno es en sí más difícil de aprender que la política, o si el rollout se distorsiona en cuanto se alarga un poco, aprender un modelo puede, al contrario, volver más frágil todo el sistema.
\end{warningbox}

\subsection{El modelo del mundo no tiene una sola forma}

Al final, el curso advierte que la forward dynamics tradicional $p(s_{t+1}\mid s_t, a_t)$ es solo un tipo de modelo del mundo. Un inverse model, un inverse model de varios pasos, un modelo que solo predice observaciones futuras e incluso la distribución conjunta de las transition tuples pueden resultar más útiles en tareas específicas.

\begin{figure}[H]
\centering
\includegraphics[width=0.9\textwidth]{slides-images/slide-09.jpg}
\caption{Las diversas formas del modelo del mundo: forward/inverse/future prediction tienen cada una su uso}
\end{figure}

\subsection{Resumen de la sección}

En esta clase, model-based RL se reubica como una \textbf{herramienta de aumento de datos}: lo importante no es lograr la simulación del mundo más perfecta, sino usar un modelo local lo bastante preciso para ampliar las muestras de entrenamiento cerca de los datos reales, manteniendo siempre el sesgo del modelo dentro de un rango aceptable.
